\documentclass[10pt,a4paper]{amsart}
\usepackage[margin=19mm]{geometry}
\usepackage{amsmath,amssymb,mathtools}
\usepackage[scr=rsfs,cal=euler]{mathalfa}
\usepackage{xfrac}
\usepackage[unicode,hidelinks,bookmarks=false]{hyperref}
\newcommand{\ie}{i.e.\ }
\newcommand{\cf}{cf.\ }
\newcommand{\resp}{respectively\ }
\newcommand{\Subsection}{\subsection}
\providecommand{\nobreakdash}{\nobreak\hbox{-}}
\setlength{\emergencystretch}{2em}
\multlinegap0pt
\usepackage{amssymb}
\newtheorem{prop}{Proposition}
\title{Schr\"odinger operators with concentric $\delta$--shell interactions}
\author{Masahiro Kaminaga}
\date{Source: arXiv:2602.09376v3 (CC BY 4.0)}
\begin{document}
\maketitle

\noindent\emph{Source and licence.} Schr\"odinger operators with concentric $\delta$--shell interactions. Version arXiv:2602.09376v3 (CC BY 4.0), distributed by arXiv under the Creative Commons Attribution 4.0 International licence, http://creativecommons.org/licenses/by/4.0/, as recorded in the arXiv licence metadata for this exact version and shown on the abstract page https://arxiv.org/abs/2602.09376v3. Full licence notice: \texttt{licenses/arxiv-2602.09376-CC-BY-4.0.txt}.\par\smallskip

\noindent\textit{Editorial scope.} Counted unit: the article's prop prop:neg-crit-corrected (source0.tex lines 1790--1810). The article's complete original proof of it is delivered with it (source0.tex lines 1813--1971, one \texttt{proof} environment of 159 lines). No mathematics is written, repaired or re-derived: the statement and the proof are the article's own lines, in the article's own notation, and no retained line is substituted or re-flowed.  Retained uncounted as the source-local context the proof needs: lines 1651--1686; lines 1666--1669.  Nothing was omitted from any retained span, so the package carries no omission record.  No retained line contains a citation, so the wrapper carries no bibliography and no bibliography entry.  No retained line contains a figure environment, an image-inclusion command or drawing code, so no image is delivered and the package declares no asset dependency.  Editorial formatting only, in addition: an AMS \texttt{amsart} preamble that restates the article's own declarations and macros used by the retained lines, namely the environments \texttt{prop} on separate counters, together with the standard packages those lines need.  The retained environments print in this document as Proposition 1 in order of appearance, whereas the article prints the same items with the numbering of its own full document.  The complete native source of the article as distributed in the arXiv e-print for this exact version is bundled unchanged as \texttt{sources/194314/source0.tex}.
\begin{prop}\label{prop:s-wave-secular}
Let $0<R_1<R_2$, $\alpha_1,\alpha_2\in\mathbb R$, and set $d=R_2-R_1$.
In the $s$--wave sector ($\ell=0$) we write $u(r)=r f(r)$.
Then $u$ satisfies
\begin{equation}\label{1D}
-\,u''(r)=E\,u(r)\quad \text{for } r\ne R_1,R_2,\qquad
u(0)=0,\qquad u\in L^2(0,\infty),
\end{equation}
together with the interface conditions
\begin{equation}\label{u-jump}
u \ \text{is continuous at } R_j,\qquad
u'(R_j+0)-u'(R_j-0)=\alpha_j\,u(R_j),\quad j=1,2.
\end{equation}
For $E=-\kappa^2<0$ the $s$--wave eigenvalues are precisely those numbers
$E=-\kappa^2$ with $\kappa>0$ which satisfy
\begin{equation}\label{eq:sec}
F_d(\kappa):=\Bigl[\gamma_1(\kappa)+(\alpha_2+\kappa)\Bigr]\cosh(\kappa d)
+\Bigl[\kappa+\gamma_1(\kappa)+\frac{\alpha_2\,\gamma_1(\kappa)}{\kappa}\Bigr]\sinh(\kappa d)=0,
\end{equation}
where
\begin{equation}\label{gamma1}
\gamma_1(\kappa)=\alpha_1+\kappa\,\coth(\kappa R_1).
\end{equation}
Equivalently, nontrivial solutions exist if and only if $\det M(\kappa;d)=0$,
where
\begin{equation}\label{eq:matrix}
M(\kappa;d)=
\left(
\begin{array}{cc}
-\gamma_1(\kappa) & \kappa\\[4pt]
(\alpha_2+\kappa)\cosh(\kappa d)+\kappa\sinh(\kappa d) &
(\alpha_2+\kappa)\sinh(\kappa d)+\kappa\cosh(\kappa d)
\end{array}
\right).
\end{equation}
\end{prop}

\begin{equation}
F_d(\kappa):=\Bigl[\gamma_1(\kappa)+(\alpha_2+\kappa)\Bigr]\cosh(\kappa d)
+\Bigl[\kappa+\gamma_1(\kappa)+\frac{\alpha_2\,\gamma_1(\kappa)}{\kappa}\Bigr]\sinh(\kappa d)=0,
\end{equation}

\begin{prop}\label{prop:neg-crit-corrected}
Let $H$ be the double $\delta$--shell Hamiltonian with radii $0<R_1<R_2$ and
couplings $\alpha_1,\alpha_2\in\mathbb R$.
In the $s$--wave subspace, the negative eigenvalues $E=-\kappa^2$ are in
one--to--one correspondence with the positive roots $\kappa>0$ of the secular
equation~\eqref{eq:sec}.

The associated $s$--wave quadratic form is given by
\[
q_0[f]
=
\int_0^\infty |f'(r)|^2\,dr
+\alpha_1|f(R_1)|^2+\alpha_2|f(R_2)|^2,
\qquad f\in H^1(0,\infty).
\]
Hence the negative part of $q_0$ depends only on the two point evaluations
$f\mapsto f(R_1)$ and $f\mapsto f(R_2)$.
As a consequence, the $s$--wave radial operator can have at most two negative
eigenvalues (counted with multiplicity), and therefore the secular
equation~\eqref{eq:sec} admits at most two positive roots.
\end{prop}

\begin{proof}
By Proposition~\ref{prop:s-wave-secular} the negative $s$--wave eigenvalues
$E=-\kappa^2$ are characterized by the positive roots $\kappa>0$ of
$F_d(\kappa)=0$, where $F_d$ is given in \eqref{eq:sec}.
We therefore study the function $F_d(\kappa)$.

The $s$--wave secular equation has the form
$$
F_d(\kappa)
=
A(\kappa)\cosh(\kappa d)+B(\kappa)\sinh(\kappa d),
\qquad \kappa>0,
$$
where
\begin{eqnarray*}
A(\kappa)&=&\gamma_1(\kappa)+\alpha_2+\kappa,\\
B(\kappa)&=&\kappa+\gamma_1(\kappa)+\frac{\alpha_2\gamma_1(\kappa)}{\kappa},\\
\gamma_1(\kappa)&=&\alpha_1+\kappa\coth(\kappa R_1).
\end{eqnarray*}
Using $\cosh t=\tfrac12(e^t+e^{-t})$ and $\sinh t=\tfrac12(e^t-e^{-t})$, we obtain
\begin{equation}\label{eq:Fd-split-new}
F_d(\kappa)=\tfrac12\bigl[F_\infty(\kappa)e^{\kappa d}+G(\kappa)e^{-\kappa d}\bigr],
\end{equation}
where
\begin{equation}\label{eq:Finfty-G}
F_\infty(\kappa)=\frac{(\kappa+\gamma_1(\kappa))(2\kappa+\alpha_2)}{\kappa},\qquad
G(\kappa)=\alpha_2\Bigl(1-\frac{\gamma_1(\kappa)}{\kappa}\Bigr).
\end{equation}

\smallskip
\noindent
We first show that the $s$--wave radial operator has at most two negative
eigenvalues.
Assume for contradiction that it has at least three negative eigenvalues and
let $L$ be the span of three corresponding eigenfunctions.
Then $\dim L=3$ and
$$
q_0[f]=(h_0f,f)_{L^2(0,\infty)}<0\qquad \text{for all }0\ne f\in L.
$$
The linear map $f\mapsto (f(R_1),f(R_2))\in\mathbb C^2$ cannot be injective on
$L$, hence there exists $0\ne f\in L$ such that $f(R_1)=f(R_2)=0$.
For this $f$ we have
$$
q_0[f]=\int_0^\infty |f'(r)|^2\,dr>0,
$$
since $\int_0^\infty |f'(r)|^2\,dr=0$ would imply that $f$ is constant, and the
conditions $f(R_1)=f(R_2)=0$ would force $f\equiv0$.
This contradiction shows $N_-(h_0)\le2$, and therefore $F_d(\kappa)=0$ has at
most two positive roots, counted with multiplicity.

\smallskip
\noindent
(i) We first consider the case $\alpha_1,\alpha_2\ge0$.
Then $\gamma_1(\kappa)\ge0$, $A(\kappa)>0$, and $B(\kappa)\ge0$ for all
$\kappa>0$.
Since $\cosh(\kappa d)>0$ and $\sinh(\kappa d)>0$ for all $\kappa>0$ and $d>0$,
we have $F_d(\kappa)>0$ for all $\kappa>0$.
Hence, there is no positive root of $F_d$, and therefore no negative $s$--wave
eigenvalue.

\smallskip
\noindent
(ii) Next we assume that exactly one shell is attractive and that the
corresponding one--shell problem supports an $s$--wave bound state.

\smallskip
\noindent
Case 1: $\alpha_1<0\le\alpha_2$ and $\alpha_1<-1/R_1$.
Then, there exists a unique $\kappa_{\mathrm{in}}>0$ satisfying
$$
\kappa+\gamma_1(\kappa)=0,
\qquad\text{that is,}\qquad
\alpha_1+\kappa\coth(\kappa R_1)+\kappa=0.
$$
At $\kappa=\kappa_{\mathrm{in}}$ one has $F_\infty(\kappa_{\mathrm{in}})=0$.
Moreover $F_\infty'(\kappa_{\mathrm{in}})\ne0$.
Indeed, $2\kappa_{\mathrm{in}}+\alpha_2>0$ and
$\kappa_{\mathrm{in}}+\gamma_1(\kappa_{\mathrm{in}})=0$ imply
$F_\infty'(\kappa_{\mathrm{in}})>0$.
Expanding \eqref{eq:Fd-split-new} near $\kappa_{\mathrm{in}}$ yields
$$
\kappa(d)=\kappa_{\mathrm{in}}
-\frac{G(\kappa_{\mathrm{in}})}{F'_\infty(\kappa_{\mathrm{in}})}e^{-2\kappa_{\mathrm{in}}d}
+O(e^{-4\kappa_{\mathrm{in}}d}),
$$
so for all sufficiently large $d$ there exists exactly one positive root of
$F_d$.
The corresponding eigenvalue gives the unique negative $s$--wave eigenvalue in
this case.

\smallskip
\noindent
Case 2: $\alpha_2<0\le\alpha_1$ and $\alpha_2<-1/R_2$.
For the single outer shell at $r=R_2$, the $s$--wave bound state condition is
$\alpha_2<-1/R_2$, and the corresponding root is given by the unique positive
zero of $2\kappa+\alpha_2$, namely
$$
\kappa_{\mathrm{out}}=-\frac{\alpha_2}{2}>0.
$$
At $\kappa=\kappa_{\mathrm{out}}$ we have $F_\infty(\kappa_{\mathrm{out}})=0$.
Moreover
$$
F_\infty'(\kappa_{\mathrm{out}})
=\frac{2\bigl(\kappa_{\mathrm{out}}+\gamma_1(\kappa_{\mathrm{out}})\bigr)}{\kappa_{\mathrm{out}}}
\ne0,
$$
and since $\alpha_1\ge0$ implies $\gamma_1(\kappa_{\mathrm{out}})>0$, we have
$F_\infty'(\kappa_{\mathrm{out}})>0$.
Expanding \eqref{eq:Fd-split-new} near $\kappa_{\mathrm{out}}$ yields
$$
\kappa(d)=\kappa_{\mathrm{out}}
-\frac{G(\kappa_{\mathrm{out}})}{F'_\infty(\kappa_{\mathrm{out}})}e^{-2\kappa_{\mathrm{out}}d}
+O(e^{-4\kappa_{\mathrm{out}}d}).
$$
Thus, we again obtain exactly one positive root for all sufficiently large $d$.
This root gives the unique negative $s$--wave eigenvalue in this case.

\smallskip
\noindent
(iii) Finally we consider the case when both shells are attractive,
$\alpha_1<0$ and $\alpha_2<0$.
Assume that each one--shell subsystem supports an $s$--wave bound state, that is,
$\alpha_j<-1/R_j$ for $j=1,2$.
Let $\kappa_{\mathrm{in}},\kappa_{\mathrm{out}}>0$ be the corresponding isolated
roots, that is, the zeros of $\kappa+\gamma_1(\kappa)$ and $2\kappa+\alpha_2$,
respectively.
Then
$$
F_\infty(\kappa_{\mathrm{in}})=0,\qquad F_\infty(\kappa_{\mathrm{out}})=0.
$$
Assume that these zeros are simple, that is,
$F_\infty'(\kappa_{\mathrm{in}})\ne0$ and $F_\infty'(\kappa_{\mathrm{out}})\ne0$.
(If one of them is not simple, then the corresponding critical tuning is treated
separately in the tunneling regime.)
For large $d$ the roots are only slightly perturbed, and we can write
\begin{eqnarray}
\delta_{\mathrm{in}}(d)&=&-\frac{G(\kappa_{\mathrm{in}})}{F'_\infty(\kappa_{\mathrm{in}})}e^{-2\kappa_{\mathrm{in}}d}
+O(e^{-4\kappa_{\mathrm{in}}d}),\nonumber\\
\delta_{\mathrm{out}}(d)&=&-\frac{G(\kappa_{\mathrm{out}})}{F'_\infty(\kappa_{\mathrm{out}})}e^{-2\kappa_{\mathrm{out}}d}
+O(e^{-4\kappa_{\mathrm{out}}d}),\nonumber
\end{eqnarray}
so for all sufficiently large $d$ there are two distinct positive roots of $F_d$,
namely $\kappa_{\mathrm{in}}+\delta_{\mathrm{in}}(d)$ and
$\kappa_{\mathrm{out}}+\delta_{\mathrm{out}}(d)$.

For general values of $d$, the number of positive roots in $(0,\infty)$ can
change only when $F_d$ has a multiple root $\kappa_0>0$, that is,
\begin{equation}\label{eq:double-root-new}
\begin{cases}
A(\kappa_0)\cosh(\kappa_0d)+B(\kappa_0)\sinh(\kappa_0d)=0,\\
(A'(\kappa_0)+dB(\kappa_0))\cosh(\kappa_0d)
+(B'(\kappa_0)+dA(\kappa_0))\sinh(\kappa_0d)=0.
\end{cases}
\end{equation}
We exclude here the threshold $\kappa=0$.
This gives a codimension--one condition on the parameters $(\alpha_1,\alpha_2,d)$.
In the generic situation the two roots persist and merge only at isolated
critical values of $(\alpha_1,\alpha_2,d)$.
\end{proof}

\end{document}
